% ===========================================================================
% Author: Adam | Module: AI Job Interview Simulator
% Project: SmartCareer AI, IA for Software Engineering 2026-2027
%
% Preambule commun au chapitre « Simulateur d'entretien ».
%
% Chargé (a) par main.tex pour la compilation autonome, (b) par le rapport
% d'équipe avant % ===========================================================================
% Author: Adam | Module: AI Job Interview Simulator
% Project: SmartCareer AI, IA for Software Engineering 2026-2027
%
% CHAPITRE « Simulateur d'entretien » — fichier à insérer dans le rapport
% d'équipe. Il ne contient aucun préambule : il enchaîne les sections.
%
% Intégration dans le rapport d'équipe (depuis la racine du dépôt) :
%
%   % --- dans le préambule du rapport d'équipe ---
%   \newcommand{\adamroot}{interview_adam/docs/report/}
%   \input{\adamroot report-preamble}
%
%   % --- dans le corps du document, à l'emplacement du chapitre ---
%   \input{\adamroot chapter-interview-adam}
%
% Si \adamroot n'est pas défini, il vaut {} (voir report-preamble.tex) : il
% faut alors compiler depuis le dossier interview_adam/docs/report/.
%
% Le fichier ne définit ni \chapter ni \section{Module 3...} : c'est au
% rapport d'équipe de donner le titre du chapitre, pour que la numérotation
% suive celle du document hôte.
%
% Références bibliographiques : le chapitre n'utilise que deux clés,
%   reimers2019sbert   (Reimers & Gurevych, Sentence-BERT, EMNLP 2019)
%   johnson2019faiss   (Johnson, Douze & Jégou, FAISS, arXiv:1702.08734)
% Elles sont définies dans la bibliographie de main.tex ; dans le rapport
% d'équipe, il faut les reprendre dans refs.bib (ou le fichier équivalent).
% Les figures sont sous \adamroot figures/ ; les fragments de code verbatim
% sont des fichiers \adamroot sections/lst-*.tex inclus par les sections.
% ===========================================================================

\input{\adamroot sections/01-introduction}
\input{\adamroot sections/02-cahier-des-charges}
\input{\adamroot sections/03-architecture}
\input{\adamroot sections/04-stack}
\input{\adamroot sections/05-modele}
\input{\adamroot sections/06-prompts}
\input{\adamroot sections/07-rag}
\input{\adamroot sections/08-securite}
\input{\adamroot sections/09-performances}
\input{\adamroot sections/10-tests}
\input{\adamroot sections/11-limites}
\input{\adamroot sections/12-conclusion}
\input{\adamroot sections/a-developpement-ia}
\input{\adamroot sections/b-demo}
\input{\adamroot sections/c-environnement}
\input{\adamroot sections/d-contribution}
. Chaque paquet est protégé
% par \@ifpackageloaded : un paquet déjà chargé ailleurs n'est pas rechargé,
% ce qui évite les « Option clash for ... ».
%
% Convention de chemin — \adamroot désigne le dossier interview_adam/docs/report/
% relativement au répertoire où l'on lance pdflatex :
%   compilation autonome (depuis docs/report/)   : \adamroot = {}   (défini ici)
%   rapport d'équipe     (depuis la racine)      : \adamroot = {interview_adam/docs/report/}
% ===========================================================================

\makeatletter

\providecommand{\adamroot}{}

% --- encodage et typographie francaise ------------------------------------
% Le chapitre doit compiler avec pdflatex (Overleaf, rapport d'equipe) MAIS
% aussi avec xelatex/lualatex : les paquets d'encodage ci-dessous sont
% specifiques a pdfTeX et font double emploi avec la gestion unicode native
% des autres moteurs. D'ou les garde-fous.
\newif\ifadam@ptex
\ifdefined\XeTeXversion\else\ifdefined\luatexversion\else\adam@ptextrue\fi\fi

\ifadam@ptex
  \@ifpackageloaded{inputenc}{}{\usepackage[utf8]{inputenc}}
  \@ifpackageloaded{fontenc}{}{\usepackage[T1]{fontenc}}
  \@ifpackageloaded{lmodern}{}{\usepackage{lmodern}}
\fi
\@ifpackageloaded{textcomp}{}{\usepackage{textcomp}}
\@ifpackageloaded{babel}{}{\usepackage[french]{babel}}
\@ifpackageloaded{microtype}{}{%
  \usepackage{microtype}%
  % L'expansion typographique de microtype n'existe qu'en pdfTeX.
  \ifadam@ptex\microtypesetup{protrusion=true,expansion=true}\fi
}

% --- boîtes, tableaux, figures, maths -------------------------------------
\@ifpackageloaded{graphicx}{}{\usepackage{graphicx}}
\@ifpackageloaded{amsmath}{}{\usepackage{amsmath}}
\@ifpackageloaded{amssymb}{}{\usepackage{amssymb}}
\@ifpackageloaded{booktabs}{}{\usepackage{booktabs}}
\@ifpackageloaded{array}{}{\usepackage{array}}
\@ifpackageloaded{tabularx}{}{\usepackage{tabularx}}
\@ifpackageloaded{longtable}{}{\usepackage{longtable}}
\@ifpackageloaded{multirow}{}{\usepackage{multirow}}
\@ifpackageloaded{xcolor}{}{\usepackage{xcolor}}
\@ifpackageloaded{caption}{}{\usepackage[font=small,labelfont=bf,skip=4pt]{caption}}
\@ifpackageloaded{float}{}{\usepackage{float}}
\@ifpackageloaded{enumitem}{}{\usepackage{enumitem}}
\@ifpackageloaded{needspace}{}{\usepackage{needspace}}

% --- dessins --------------------------------------------------------------
\@ifpackageloaded{tikz}{}{\usepackage{tikz}}
\usetikzlibrary{arrows.meta,positioning,shapes.geometric,calc,fit,backgrounds}

% --- code et verbatim -----------------------------------------------------
\@ifpackageloaded{listings}{}{\usepackage{listings}}

% --- géométrie : seulement si le document ne l'a pas déjà fixée -----------
\@ifpackageloaded{geometry}{}{\usepackage[a4paper,top=2.4cm,bottom=2.4cm,left=2.6cm,right=2.6cm]{geometry}}

% --- caractères unicode hors T1 (utilisés dans les verbatims de prompts) --
\@ifpackageloaded{newunicodechar}{}{\usepackage{newunicodechar}}
\newunicodechar{→}{\ensuremath{\rightarrow}}
\newunicodechar{←}{\ensuremath{\leftarrow}}
\newunicodechar{≈}{\ensuremath{\approx}}
\newunicodechar{≥}{\ensuremath{\ge}}
\newunicodechar{≤}{\ensuremath{\le}}
\newunicodechar{≠}{\ensuremath{\neq}}
\newunicodechar{…}{\ldots}
\newunicodechar{—}{\textemdash}
\newunicodechar{–}{\textendash}
\newunicodechar{°}{\textdegree}
\newunicodechar{✓}{\ensuremath{\surd}}

% --- hyperref en dernier --------------------------------------------------
\@ifpackageloaded{hyperref}{}{\usepackage[hidelinks]{hyperref}}

\makeatother

% --- palette --------------------------------------------------------------
\providecolor{adamInk}{HTML}{1F2933}
\providecolor{adamAccent}{HTML}{8E1F3F}
\providecolor{adamSoft}{HTML}{F4F1EE}
\providecolor{adamKey}{HTML}{1A5FB4}
\providecolor{adamStr}{HTML}{0B6E4F}
\providecolor{adamCom}{HTML}{7A6A55}
\providecolor{adamFrame}{HTML}{C9B8AE}

% --- macros du chapitre ---------------------------------------------------
% Cellule de tableau : alignee a gauche et autorisant la coupure des
% identificateurs longs (un raggedright classique interdit toute césure, ce
% qui fait deborder les colonnes etroites).
\makeatletter
\@ifundefined{cellflush}{%
  \newcommand{\cellflush}{%
    \raggedright\arraybackslash
    \hyphenpenalty=50\relax
    \exhyphenpenalty=100\relax}%
}{}
% Le tirete bas souligne devient un point de coupure : les noms
% MAX_CONCURRENT_LLM_CALLS ou interview_adam/... tiennent dans les colonnes.
\let\_\relax
\DeclareRobustCommand{\_}{\textunderscore\allowbreak}
\makeatother
% Identificateur en police clavier, cassable apres chaque underscore : les
% noms longs (MAX_CONCURRENT_LLM_CALLS, interview_adam/...) tiennent dans les
% colonnes etroites des tableaux.
\makeatletter
\providecommand{\adamtt}[1]{{\def\_{\textunderscore\allowbreak}\ttfamily #1}}
\makeatother
\providecommand{\fig}[1]{\adamroot figures/#1}
\providecommand{\file}[1]{\adamtt{#1}}
\providecommand{\code}[1]{\adamtt{#1}}
\providecommand{\const}[1]{\adamtt{#1}}
\providecommand{\route}[1]{\adamtt{#1}}
\providecommand{\var}[1]{\textit{#1}}
\providecommand{\NM}{\textup{\textsc{non mesuré}}}
\providecommand{\note}[1]{\textcolor{adamAccent}{#1}}
% Capture d'ecran de demonstration : largeur plafonnee, hauteur plafonnee pour
% que les captures verticales (defilement complet) tiennent sur une page.
\providecommand{\shot}[1]{\includegraphics[%
  width=0.98\textwidth,totalheight=0.66\textheight,keepaspectratio]{\fig{#1}}}

% --- styles listings ------------------------------------------------------
\lstdefinestyle{adambase}{%
  basicstyle=\ttfamily\scriptsize,
  keywordstyle=\color{adamKey}\bfseries,
  commentstyle=\color{adamCom}\itshape,
  stringstyle=\color{adamStr},
  showstringspaces=false,
  columns=fullflexible,
  keepspaces=true,
  breaklines=true,
  breakautoindent=true,
  postbreak=\mbox{\textcolor{adamAccent}{$\hookrightarrow$}\space},
  tabsize=4,
  frame=single,
  rulecolor=\color{adamFrame},
  backgroundcolor=\color{adamSoft},
  framesep=4pt,
  xleftmargin=6pt,
  xrightmargin=4pt,
  aboveskip=8pt,
  belowskip=8pt,
  numbers=none,
  inputencoding=utf8,
  escapeinside={(*}{*)}}

\lstdefinestyle{adamprompt}{style=adambase}
\lstdefinestyle{adamjson}{style=adambase, morekeywords={true,false,null}}
\lstdefinestyle{adamcode}{style=adambase, language=Python}
\lstdefinestyle{adamshell}{style=adambase}

\lstset{style=adambase}

% --- styles TikZ ----------------------------------------------------------
\tikzset{
  adambox/.style={draw=adamFrame, fill=white, rounded corners=2pt, align=center,
                  font=\footnotesize\sffamily, inner sep=5pt, minimum width=2.6cm,
                  text=adamInk, thick},
  adamsvc/.style={adambox, fill=adamSoft, draw=adamAccent},
  adamext/.style={adambox, dashed},
  adamdb/.style={adambox, fill=white, rounded corners=6pt, minimum width=3cm},
  adamflow/.style={->, thick, draw=adamInk, >=Stealth},
  adamback/.style={->, thick, draw=adamAccent, >=Stealth, dashed},
  lam/.style={font=\scriptsize\sffamily\itshape, text=adamCom},
  ll/.style={draw=adamFrame, fill=adamSoft, rounded corners=2pt, align=center,
             font=\scriptsize\sffamily, inner sep=4pt, text=adamInk, thick},
  msg/.style={->, thick, draw=adamInk, >=Stealth},
  rmsg/.style={->, thick, draw=adamAccent, >=Stealth, dashed},
  lifeline/.style={dotted, draw=adamFrame, thick},
  mlab/.style={font=\scriptsize\sffamily, text=adamInk, inner sep=2pt},
}

% --- tableaux -------------------------------------------------------------
% Les colonnes de largeur fixe sont écrites en toutes lettres dans les tableaux
% (>{\raggedright\arraybackslash}p{...}) : aucun \newcolumntype, donc aucune
% collision possible avec le préambule du rapport d'équipe.
\renewcommand{\arraystretch}{1.2}

% --- tolérance de mise en page --------------------------------------------
% Le chapitre est dense en identifiants longs (noms de fichiers, de constantes)
% que la coupure automatique ne sait pas découper : on laisse TeX étirer la
% ligne avant de déclarer un débordement, et on offre un point de coupure
% supplémentaire (\brk) pour les tokens typographiques sans espace.
\setlength{\emergencystretch}{3em}
\providecommand{\brk}{\allowbreak}
